\documentclass[12pt]{article}
\usepackage{amsmath}
\usepackage{amssymb}
\usepackage{amsthm}
\usepackage{cancel}
\usepackage{enumitem}
\usepackage{esdiff}
\usepackage{graphicx}
\usepackage{multicol}
\usepackage{siunitx}
\usepackage{ulem}
% \usepackage{pgfplots}
\usepackage{wrapfig}
\usepackage{xcolor}

\newcommand{\e}[1]{e^{#1}}
\newcommand{\ei}[1]{e^{i\left(#1\right)}}
\newcommand{\E}[1]{\times 10^{#1}}
\newcommand{\tagref}[1]{\tag{\ref{#1}}}

\renewcommand\qedsymbol{TENA}

\title{
    Homework \#10
    \\  \small
    PHYS 4D: Modern Physics
    \\
    Problems from Krane's textbook
    }
\author{Donald Aingworth IV}
\date{March 23, 2026}

\begin{document}
    \maketitle

    \section{Optional}
        \subsection{Chapter 7, Question 1}
            How does the quantum-mechanical interpretation of the hydrogen atom differ from the Bohr model?

            \subsubsection{Answer}

        \subsection{Chapter 7, Question 3}
            What are the meanings of the quantum numbers n, l, m l according to (a) the quantum-mechanical calculation; (b) the vector model; (c) the Bohr (orbital) model?

            \subsubsection{Answer}
            
        \subsection{Chapter 7, Question 5}
            How does the orbital angular momentum differ between the Bohr model and the quantum-mechanical calculation?

            \subsubsection{Answer}
            
        \subsection{Chapter 7, Question 11}
            The probability density $\psi* \psi$ does not depend on $\phi$ for the wave functions listed in Table 7.1. 
            What is the significance of this?

            \subsubsection{Answer}
            
        \subsection{Chapter 7, Question 13}
            Can a hydrogen atom in its ground state absorb a photon (of the proper energy) and end up in the 3d state?

            \subsubsection{Answer}
            
        \subsection{Chapter 7, Question 15}
            The photon has a spin quantum number of 1, but its spin magnetic moment is zero. 
            Explain.

            \subsubsection{Answer}
            
        \subsection{Chapter 7, Question 19}
            How would (a) the Zeeman effect and (b) the fine structure be different in a muonic hydrogen atom? 
            (See Problems 41 and 42 in Chapter 6.) 
            The muon has the same spin as the electron, but is 207 times as massive.

            \subsubsection{Answer}
            
        \subsection{Chapter 7, Problem 21}
            List the excited states (in spectroscopic notation) to which the $4p$ state can make downward transitions.

            \subsubsection{Answer}
            
        \subsection{Chapter 7, Problem 22}
            (a) A hydrogen atom is in an excited $5g$ state, from which it makes a series of transitions by emitting photons, ending in the $1s$ state. 
            Show, on a diagram similar to Figure 7.19, the sequence of transitions that can occur. 
            (b) Repeat part (a) if the atom begins in the $5d$ state.

            \subsubsection{Answer}
            
        \subsection{Chapter 7, Problem 25}
            A collection of hydrogen atoms is placed in a magnetic field of 3.50 T. 
            Ignoring the effects of electron spin, find the wavelengths of the three normal Zeeman components (a) of the 3d to 2p transition; (b) of the 3s to 2p transition.

            \subsubsection{Answer}
            
        \subsection{Chapter 7, Problem 27}
            Calculate the energies and wavelengths of the 3d to 2p transition, taking into account the fine structure of both levels.
            How many component wavelengths might there be in the transition?

            \subsubsection{Answer}
            
        \subsection{Chapter 7, Problem 34}
            For the 1s, 2s, and 2p states of hydrogen, show that $(r^{-1})_{\rm av} = 1/n^2a_0$. 
            This turns out to be a general result for any state of hydrogen. 
            Based on this result, explain why the Bohr model gives such a good estimate for the fine-structure splitting as well as for other magnetic effects due to the circulating electron.

            \subsubsection{Answer}
            
    \pagebreak
    
    \section{Chapter 8, Question 2}
        Why do the $4s$ and $3d$ subshells appear so close in energy, when they belong to different principal quantum numbers $n$?

        \subsection{Solution}
    \pagebreak
    
    \section{Chapter 8, Question 3}
        Would you expect element 107 to be a good conductor or a poor conductor? 
        How about element 111? 
        Do you expect element 112 to be paramagnetic or diamagnetic?

        \subsection{Solution}
    \pagebreak

    \section{Chapter 8, Question 10}
        Suppose we do a Stern-Gerlach experiment using an atom that has angular momentum quantum numbers L and S in its ground state. 
        Into how many components will the beam split? 
        Do you expect them to be equally spaced?

        \subsection{Solution}
    \pagebreak
    
    \section{Chapter 8, Question 17}
        The first excited state in sodium is a fine-structure doublet; the wavelengths emitted in the decay of these states are 589.59 nm and 589.00 nm, a difference of 0.59 nm.
        The excited $4s^1$ state in sodium (see Figure 8.4) decays to the $3p$ doublet with the emission of radiation at the wavelengths 1138.15 nm and 1140.38 nm, a difference of 2.23 nm. Explain how the $3p$ fine structure can give a wavelength difference of 0.59 nm in one case and 2.23 nm in the other case.

        \subsection{Solution}
    \pagebreak
    
    \section{Chapter 8, Question 20}
        Explain what is meant by a population inversion and why it is necessary for the operation of a laser.

        \subsection{Solution}
    \pagebreak
    
    \section{Chapter 8, Problem 1}
        (a) List the six possible sets of quantum numbers (n, l, ml, ms) of a 2p electron. 
        (b) Suppose we have an atom such as carbon, which has two 2p electrons. 
        Ignoring the Pauli principle, how many different possible combinations of quantum numbers of the two electrons are there? 
        (c) How many of the possible combinations of part (b) are eliminated by applying the Pauli principle? 
        (d) Suppose carbon is in an excited state with configuration $2p^1 3p^1$. 
        Does the Pauli principle restrict the choice of quantum numbers for the electrons? 
        How many different sets of quantum numbers are possible for the two electrons?
        \subsection{Solution}
    \pagebreak
    
    \section{Chapter 8, Problem 19}
        A small helium-neon laser produces a light beam with an average power of 3.5 mW and a diameter of 2.4 mm.
        (a) How many photons per second are emitted by the laser?
        (b) What is the amplitude of the electric field of the light wave? Compare this result with the electric field at a distance of 1 m from an incandescent light bulb that emits 100 W of visible light.
        \subsection{Solution}
    \pagebreak
    
    \section{Chapter 8, Problem 21}
        (a) The ionization energy of sodium is 5.14 eV. 
        What is the effective charge seen by the outer electron? 
        (b) If the $3s$ electron of a sodium atom is moved to the $4f$ state, the measured binding energy is 0.85 eV. 
        What is the effective charge seen by an electron in this state?
        \subsection{Solution}
    
\end{document}